% Part: turing-machines 
% Chapter: machines-computations 
% Section: introduction

\documentclass[../../../include/open-logic-section]{subfiles}

\begin{document}

\olfileid{tur}{mac}{int}
\olsection{تعارف}

اک فنکشن، مثلاً $\Nat$ توں $\Nat$ تک،
\emph{قابلِ حساب} ہون دا کی مطلب اے؟
ابتدائی جواباں وچوں اک، تے سب توں معروف
جواب، ایہ اے کہ جے اک ٹیورنگ مشین
اوس فنکشن دا حساب کر سکدی اے تاں
اوہ قابلِ حساب اے۔ ایہ تصور ایلن ٹیورنگ
(Alan Turing) نے 1936 وچ پیش کیتا۔
ٹیورنگ مشیناں \emph{حساب دے اک نمونے}
دی مثال نیں: ایہ «حساب کرن دے طریقے»
دے خیال نوں ریاضیاتی طور اُتے ٹھیک ٹھیک
تعریف کرن دا اک طریقہ نیں۔ اوس خیال
دا بالکل کی مطلب اے، ایس اُتے بحث
ہوندی اے؛ پر عام اتفاق اے کہ ٹیورنگ
مشیناں حساب کرن دے طریقے مقرر کرن
دا اک طریقہ نیں۔ بھاویں «ٹیورنگ مشین»
سن کے حرکت کرن والے پرزیاں دی مادی
مشین دا خیال آندا اے، پر ٹھیک ٹھیک آکھیے
تاں ٹیورنگ مشین خالص ریاضیاتی بناوٹ اے۔
ایہ حساب کرن دے طریقے دا اوہ ریاضیاتی روپ اے
جیہدے اُتے اصلی مشین دیاں عملی بندشاں نہیں لگدیاں۔ مثال لئی،
اسی ٹیورنگ مشین دے وقت یا یادداشت
دی ضرورت اُتے کوئی حد نہیں رکھدے:
اوہ ایداں دا حساب وی کر سکدی اے
جیدے لئی کائنات دے ایٹماں دی گنتی
توں ودھ ذخیرہ کرن دی تھاں یا قدم
درکار ہون۔

\begin{explain}
شاید ٹیورنگ مشین نوں اک خاص قسم
دے فرضی آلے دا پروگرام سمجھنا
سب توں چنگا اے۔ ایس آلے وچ اک
\emph{پٹی} (tape) تے اک
\emph{پڑھن-لکھن والا سر} (head)
ہوندا اے۔ ساڈی ٹیورنگ مشین دی
صورت وچ پٹی اک پاسے، یعنی سجّے
پاسے، ول لامتناہی اے۔ ایہ
\emph{خانیاں} وچ ونڈی ہوئی اے؛
ہر خانے وچ کوئی نشان ہو سکدا اے
جیڑا اک محدود \emph{حرفاں دے مجموعے}
(alphabet) وچوں لیا جاندا اے۔
ایس مجموعے وچ وکھو وکھرے
نشان کوئی وی محدود گنتی وچ
رکھ سکدے آں،
پر زیادہ تر ساڈے لئی تِن ہی کافی
ہون گے: $\TMendtape$،
$\TMblank$ تے $\TMstroke$۔
جدوں آلہ شروع ہوندا اے،
پٹی خالی ہوندی اے، یعنی
ہر خانے وچ $\TMblank$ ہوندا اے؛
بس سب توں کھبے خانے وچ
$\TMendtape$ ہوندا اے تے
محدود گنتی دے کجھ خانے
\emph{انپٹ} رکھدے نیں۔
ہر ویلے آلہ محدود گنتی
دیاں \emph{حالتاں} وچوں
کسے اک حالت وچ ہوندا اے۔
شروع وچ سر سب توں کھبے
خانے نوں پڑھدا اے تے آلہ
اک مقرر کیتی \emph{ابتدائی حالت}
وچ ہوندا اے۔ آلے دے چلن
دے ہر قدم اُتے، سر ہِٹھ
والے خانے دا نشان، آلے
دی موجودہ حالت تے ٹیورنگ
مشین دا پروگرام رل کے
اگلا عمل طے کردے نیں۔
ٹیورنگ مشین دا پروگرام
اک جزوی فنکشن نال دتا جاندا
اے: ایہ اک حالت~$q$
تے اک نشان~$\sigma$
انپٹ لے کے اک تِکا
$\tuple{q', \sigma',
  D}$ واپس دیندا اے۔
جدوں آلہ حالت $q$
وچ ہووے تے نشان $\sigma$
پڑھے، تاں اوہ موجودہ
خانے دا نشان $\sigma'$
نال بدلدا اے؛ $D$
دی قدر $\TMleft$،
$\TMright$ یا $\TMstay$
ہون مطابق سر کھبے ول سرکدا اے،
سجّے ول سرکدا اے یا اوتھے ای رہندا اے؛
تے آلہ حالت~$q'$
وچ چلا جاندا اے۔

مثال لئی، \olref{fig:tm}
وچ دکھائی حالت ویکھو۔
\begin{figure}
  \olasset{\olpath/assets/diagrams/turing-machine.tikz}
  \caption{اک ٹیورنگ مشین اپنا پروگرام چلا رہی اے۔}
  \ollabel{fig:tm}
\end{figure}
ٹیورنگ مشین دی پٹی دے
دکھائی دیندے حصے وچ
سب توں کھبے خانے اُتے
پٹی دے اختتام دا نشان
$\TMendtape$ اے؛ اوس دے
بعد تِن $1$، اک $0$
تے ہور چار $1$ نیں۔
سر کھبے پاسے توں تیجے
خانے نوں پڑھ ریا اے؛
اوتھے اک~$1$ اے تے
آلہ حالت~$q_1$ وچ اے۔
اسی آکھدے آں، «مشین
اک $1$ پڑھ رہی اے
جدوں اوہ حالت~$q_1$
وچ اے۔» جے ٹیورنگ
مشین دا پروگرام انپٹ
$\tuple{q_1, 1}$ لئی
تِکا $\tuple{q_2,
0, \TMstay}$ واپس دیوے،
تاں مشین تیجے خانے
والے~$1$ نوں اک~$0$
نال بدل دیوے گی،
پڑھن-لکھن والے سر نوں
اوتھے ہی رکھے گی تے
حالت~$q_2$ وچ چلی
جاوے گی۔ فیر جے پروگرام
تِکا $\tuple{q_3, 0, \TMright}$،
انپٹ $\tuple{q_2, 0}$
لئی واپس دیوے، تاں مشین
موجودہ~$0$ اُتے فیر
اک~$0$ لکھے گی
(یعنی سر ہِٹھ پٹی دا نشان
اوہو رہے گا)، سر نوں اک خانہ
سجّے ول سرکا دے گی تے حالت~$q_3$
وچ آ جاوے گی۔ ایس
طرح اگے چلدا اے۔

اسی آکھدے آں کہ مشین
\emph{رُک جاندی اے} جدوں
اوہ کسے حالت $q_n$ وچ
ایہو جیہا نشان $\sigma$ پڑھے جیہدے لئی
$\tuple{q_n, \sigma}$
دا کوئی حکم موجود نہ ہووے؛
یعنی انپٹ
$\tuple{q_n,\sigma}$
لئی حالتی تبدیلی
دا فنکشن غیر تعریف شدہ
ہووے۔ دوجے لفظاں وچ،
مشین کول کرن لئی
کوئی حکم نہیں رہندا،
تے ایس ویلے اوہ
چلنا بند کر دیندی اے۔
کدی کدی رُکّن نوں
اک خاص رُکّن والی
حالت~$h$ نال وی
دکھایا جاندا اے۔
ایہ اگے ودھ تفصیل
نال وکھایا جاوے گا۔
\end{explain}

\begin{digress}
ٹیورنگ دے مقالے
«قابلِ حساب اعداد بارے» (On computable numbers)
دی خوبی ایہ اے کہ
اوہ صرف رسمی تعریف
نہیں دیندا، بلکہ ایہ
دلیل وی دیندا اے
کہ تعریف قابلِ حساب ہون دے
بدیہی تصور نوں بیان کردی اے۔ تعریف توں
صاف ہونا چاہیدا اے
کہ ٹیورنگ مشین نال
قابلِ حساب ہر فنکشن
بدیہی معنی وچ قابلِ
حساب اے۔ ٹیورنگ
ایہ اُلٹی گل لئی،
یعنی جیہڑا وی فنکشن
اسی فطری طور اُتے
قابلِ حساب سمجھدے
آں اوہ وی ایداں
دی مشین نال قابلِ
حساب اے، تِن طرح
دیاں دلیلاں دیندا
اے۔ اوس دے اپنے
لفظاں وچ اوہ ایہ نیں:
\begin{enumerate}
\item بدیہی سمجھ نوں سدھا بنیاد بنانا۔
\item دو تعریفاں دی ہم ارزیت
  دا ثبوت (جے نویں تعریف
  بدیہی سمجھ نال زیادہ میل کھاندی ہووے)۔
\item قابلِ حساب اعداد دیاں وڈیاں جماعتاں دیاں مثالاں دینا۔
\end{enumerate}
ساڈا مقصد قابلِ حساب ہون
دے تصور نوں «اصولی
طور اُتے» تعریف کرن
دی کوشش اے، یعنی
وقت تے تھاں دیاں
عملی حداں نوں لحاظ
وچ لیائے بغیر۔ بے شک،
قابلِ حساب ہون دی سب توں وسیع
تعریف مقرر ہو جائے تاں محدود
وسیلیاں نال حساب نوں
وی ویکھ سکدے آں؛
ایہ «حسابی پیچیدگی»
(computational complexity)
نوں کہے جاندے
مضمون دا مرکز اے۔
\end{digress}

\begin{history}
ایلن ٹیورنگ (Alan Turing)
نے 1936 وچ ٹیورنگ
مشیناں ایجاد کیتیاں۔
اوس ویلے اوس دی
دلچسپی پہلے درجے دی
منطق دے قابلِ فیصلہ
ہون وچ سی، پر ایس
مقالے نوں کمپیوٹر
ڈیزائن دیاں بنیاداں
بابت فیصلہ کن مقالہ
وی آکھیا گیا اے۔
مقالے وچ ٹیورنگ دی
توجہ قابلِ حساب
حقیقی اعداد اُتے
اے، یعنی اوہ حقیقی
اعداد جِنہاں دے
اعشاری پھیلاؤ قابلِ
حساب نیں؛ پر اوہ
ایہ وی دسدا اے
کہ اپنے تصورات نوں
قدرتی اعداد اُتے
قابلِ حساب فنکشناں
وغیرہ لئی ڈھالنا
اوکھا نہیں۔ یاد رکھو
کہ ایہ پہلا کم کرن
والا عام مقصد دا
کمپیوٹر 1941 وچ
بنن توں پورے پنج
سال پہلے دی گل اے؛
اوہ کمپیوٹر جرمن
کونراڈ زُوزے (Konrad Zuse)
نے اپنے ماں پیو دے گھر دی
بیٹھک وچ بنایا سی۔ ایہ
ٹیورنگ تے اوس دے
بلیچلے پارک (Bletchley Park)
دے ساتھیواں ولوں
کوڈ توڑن والا کولوسس
(Colossus) بناؤن
(1943) توں ست سال،
امریکی ENIAC (1945)
توں نو سال، مانچسٹر
وچ پہلا برطانوی
عام مقصد دا کمپیوٹر
— Manchester Small-Scale
Experimental Machine —
بنن (1948) توں باراں
سال، تے امریکیواں
ولوں BINAC دی پہلی
آزمائش (1949) توں
تیرہ سال پہلے سی۔
مانچسٹر SSEM نوں
پہلا ایسا کمپیوٹر
ہون دا امتیاز حاصل
اے جیڑا پروگرام
اپنی یادداشت وچ
رکھدا سی؛ ایس توں
پہلیاں مشیناں نوں
ہر نویں کم لئی
ہتھ نال تارا دوبارہ
جوڑݨے پیندے سن۔
\end{history}

\end{document}
